\documentclass[11pt]{article}
\usepackage[margin=1in]{geometry}
\usepackage[T1]{fontenc}
\usepackage[utf8]{inputenc}
\usepackage{lmodern}
\usepackage{xcolor}
\usepackage{listings}
\usepackage{hyperref}
\lstset{
  basicstyle=\ttfamily\small,
  columns=fullflexible,
  keepspaces=true,
  frame=single,
  breaklines=true,
  showstringspaces=false
}
\title{pure0: A Tiny Concatenative Language for Sound Synthesis}
\author{}
\date{}
\begin{document}
\maketitle

\section*{Overview}
\texttt{pure0} is a very small Forth-like language for sound synthesis. It is stack-based and concatenative: words consume values from the stack and push results back. The central idea is that numbers, strings, and signals can all live on the stack. A block of the form \texttt{\{ ... \}} builds a signal by evaluating its body sample by sample.

\section*{Core ideas}
\begin{itemize}
\item Numbers are plain scalar values.
\item Strings are written in double quotes.
\item Signals are vectors of samples, possibly multichannel.
\item Words operate by popping from the stack and pushing results.
\item \texttt{: name ... ;} defines a new word.
\end{itemize}

For example:
\begin{lstlisting}
: sr 48000 ;
: tone440 0.2 2 { 440 } osc ;
\end{lstlisting}

Now \texttt{sr} pushes \texttt{48000} and \texttt{tone440} synthesizes a 2-second sine wave at 440 Hz.

\section*{What \texttt{\{ ... \}} means}
A block is a signal generator. The two values before the block are interpreted as:
\begin{itemize}
\item amplitude context \texttt{amp}
\item duration in seconds
\end{itemize}

So:
\begin{lstlisting}
0.2 2 { 440 }
\end{lstlisting}
means: generate a 2-second signal by evaluating the body once per sample, with local variables:
\begin{itemize}
\item \texttt{amp}: the amplitude value passed before the block
\item \texttt{t}: current time in seconds
\item \texttt{ph}: normalized phase from 0 to 1 across the block
\end{itemize}

In this case the body always returns \texttt{440}, so the result is a constant control signal equal to 440. This is useful as a frequency signal for \texttt{osc}.

\section*{Built-in variables}
By default:
\begin{lstlisting}
sr   \textit{sample rate, default 44100}
ch   \textit{number of channels, default 1}
\end{lstlisting}

You can inspect variables and words with:
\begin{lstlisting}
.v
\end{lstlisting}

You can define a numeric variable dynamically with:
\begin{lstlisting}
"sr" 48000 def
\end{lstlisting}

Or more idiomatically define a word:
\begin{lstlisting}
: sr 48000 ;
\end{lstlisting}

In the current interpreter, the DSP engine reads the numeric variable table for \texttt{sr} and \texttt{ch}. So if you want to change the actual rendering sample rate or channel count, use \texttt{def}:
\begin{lstlisting}
"sr" 48000 def
"ch" 2 def
\end{lstlisting}

\section*{Stack tools}
\begin{lstlisting}
.s    \textit{print full stack}
.     \textit{print and pop top item}
dup   \textit{duplicate top}
swap  \textit{swap top two}
drop  \textit{drop top item}
\end{lstlisting}

Example:
\begin{lstlisting}
1 2 3 .s
2 3 + .
\end{lstlisting}

\section*{Arithmetic}
The words \texttt{+ - * /} work on numbers, and also pointwise on signals. Strings can be concatenated with \texttt{+}.

\begin{lstlisting}
2 3 + .
"out" ".wav" + .
\end{lstlisting}

\section*{Oscillators and envelopes}
\subsection*{Oscillator}
\texttt{osc} now expects a frequency signal and returns a sine wave. It does not apply amplitude internally.

A minimal oscillator:
\begin{lstlisting}
0.2 2 { 440 } osc
\end{lstlisting}

This creates a 2-second sine wave at 440 Hz.

\subsection*{Envelopes}
\begin{lstlisting}
env-    \textit{constant envelope}
env<    \textit{fade in}
env>    \textit{fade out}
env<>   \textit{fade in + fade out}
adsr    \textit{ADSR envelope}
bpf     \textit{breakpoint envelope}
\end{lstlisting}

Examples:
\begin{lstlisting}
0.3 2 env-
0.3 2 env<>
0.3 2 0.01 0.2 0.6 0.3 adsr
\end{lstlisting}

To apply an envelope to a signal, use \texttt{vmul}:
\begin{lstlisting}
0.3 2 env<>
0.2 2 { 440 } osc
vmul
\end{lstlisting}

\section*{Breakpoint functions}
\texttt{bpf} takes pairs \texttt{x y}, then the number of points, then a duration:
\begin{lstlisting}
0 0  0.1 1  0.7 0.4  1 0  4 2 bpf
\end{lstlisting}
This means 4 points over 2 seconds. \texttt{x} should normally go from 0 to 1.

It is useful as an envelope:
\begin{lstlisting}
0 0  0.1 1  0.7 0.4  1 0  4 2 bpf
0.2 2 { 440 } osc
vmul
\end{lstlisting}

\section*{Signal combinators}
\begin{lstlisting}
vmul    \textit{pointwise multiply}
vadd    \textit{pointwise add}
mix     \textit{sum N signals}
concat  \textit{append signals in time}
delay   \textit{prepend silence}
stereo  \textit{combine two mono signals into stereo}
\end{lstlisting}

Examples:
\subsection*{Mix two tones}
\begin{lstlisting}
0.2 2 { 440 } osc
0.2 2 { 660 } osc
2 mix
\end{lstlisting}

\subsection*{Concatenate two tones}
\begin{lstlisting}
0.2 1 { 440 } osc
0.2 1 { 660 } osc
concat
\end{lstlisting}

\subsection*{Delay a tone}
\begin{lstlisting}
0.2 1 { 440 } osc
0.5 delay
\end{lstlisting}

\subsection*{Stereo output}
\begin{lstlisting}
"ch" 2 def
0.2 2 { 440 } osc
0.2 2 { 660 } osc
stereo
\end{lstlisting}

\section*{Strings, file names, and wavread}
Strings are written in double quotes:
\begin{lstlisting}
"out.wav"
\end{lstlisting}

Read a wav file:
\begin{lstlisting}
"input.wav" wavread
\end{lstlisting}

Render with an explicit output name:
\begin{lstlisting}
0.3 2 env<>
0.2 2 { 440 } osc
vmul
"note.wav" render
\end{lstlisting}

If you omit the name, \texttt{render} writes to \texttt{out.wav}.

\section*{Typical examples}
\subsection*{A simple note}
\begin{lstlisting}
0.3 2 env<>
0.2 2 { 440 } osc
vmul
"note.wav" render
\end{lstlisting}

\subsection*{A two-note phrase}
\begin{lstlisting}
0.3 1 env<>
0.2 1 { 440 } osc
vmul

0.3 1 env<>
0.2 1 { 660 } osc
vmul

concat
"phrase.wav" render
\end{lstlisting}

\subsection*{A gliding tone}
\begin{lstlisting}
0.3 3 env<>
0.2 3 { 220 220 ph * + } osc
vmul
"glide.wav" render
\end{lstlisting}

\subsection*{An ADSR note}
\begin{lstlisting}
0.3 2 0.01 0.2 0.6 0.3 adsr
0.2 2 { 330 } osc
vmul
"adsr.wav" render
\end{lstlisting}

\subsection*{Stereo example}
\begin{lstlisting}
"ch" 2 def
0.25 2 env<>
0.2 2 { 440 } osc
vmul

0.25 2 env<>
0.2 2 { 660 } osc
vmul

stereo
"stereo.wav" render
\end{lstlisting}

\section*{REPL and files}
If you launch \texttt{pure0} with no files, it starts a REPL using \texttt{getline}. This is convenient for a concatenative language.

If you pass a file, it executes the file contents.

\section*{Remarks}
This is a toy language. The current interpreter is intentionally small and conservative. The core aesthetic is:
\begin{itemize}
\item keep syntax minimal
\item keep the stack visible
\item treat blocks as pure signal builders
\item compose by concatenation of words
\end{itemize}

A natural next step would be to add proper scope for local definitions, better multichannel signal constructors, and richer file I/O.

\end{document}
